\section{Method}
\label{sec:method}

The final controller has two components: a frozen transit-adapted H2O+ actor
and a deterministic physical-action transform. Learning supplies a continuous
state-dependent holding proposal; the transform calibrates its magnitude and
applies a bounded service-completion reserve. No parameter is updated during
target-SUMO evaluation.

\subsection{Structured hurdle actor}
\label{sec:method:architecture}

The actor, critics, and value function use separate structured state encoders
with the same topology. A shared slot network encodes each of the 16 possible
event slots. Five pooled views---the current event, active-slot mean,
active-slot maximum, previously selected-slot mean, and future-slot
mean---are concatenated with the five-value batch prefix and mapped to a
256-dimensional representation. This preserves the ordered action prefix of a
simultaneous batch while masking inactive padding.

Because exact zero hold is frequent and operationally distinct from a small
positive duration, the actor is a hurdle distribution. Let
$p^+_{\phi}(s)=\sigmoid(g_{\phi}(s))$ be the positive-hold
probability, where $g_{\phi}(s)$ is the gate logit and
$\sigmoid(\zeta)=1/(1+\exp(-\zeta))$. Conditional on holding, the latent variable
$\zeta\sim\mathcal{N}(\mu_{\phi}(s),\sigma_{\phi}^2(s))$ follows a normal
distribution with mean $\mu_{\phi}(s)$ and positive standard deviation
$\sigma_{\phi}(s)$. Its positive-branch action is
$a^+=2\sigmoid(\zeta)-1\in(-1,1)$. The policy measure is
\begin{equation}
\pi_{\phi}(\mathrm{d}a\mid s)=
\left(1-p^+_{\phi}(s)\right)\delta_{-1}(\mathrm{d}a)
+p^+_{\phi}(s)f_{\phi}(a\mid s)
\mathbb{I}\{a\in(-1,1)\}\,\mathrm{d}a,
\label{eq:hurdle_policy}
\end{equation}
where $\delta_{-1}$ is the unit point mass at the no-hold action,
$f_{\phi}(a\mid s)$ is the conditional logistic-normal action density,
$\mathbb{I}\{\cdot\}$ is one when its condition holds and zero otherwise,
and $\mathrm{d}a$ denotes the continuous-action measure. For likelihoods and
entropy, we use the density $\varpi_{\phi}(a\mid s)$ with respect to the common
point-mass-plus-Lebesgue measure: it is $1-p^+_{\phi}(s)$ at $a=-1$ and
$p^+_{\phi}(s)f_{\phi}(a\mid s)$ for $-1<a<1$.
\Cref{prop:hurdle} in \Cref{app:properties:hurdle} derives the density and
normalization, and shows that $2\sigmoid(\mu_{\phi}(s))-1$ is the positive
branch's median, not generally its mean.

Deterministic evaluation uses that median when $p^+_{\phi}(s)\geq0.5$ and
otherwise returns $a=-1$. Physical holds are obtained through
\Cref{eq:v7_action} and remain in $[0,60]$ s. The implementation keeps the
continuous branch strictly inside $(-1,1)$ by clipping at floating-point
precision; the operational transform can still apply the exact 60-s cap.

The critic has two independently parameterized heads, with two target heads
updated by Polyak averaging at rate 0.005. The actor, critic, and value heads
each use two 256-unit hidden layers after the structured encoder. Further
implementation values are given in \Cref{app:architecture,app:training}.

\subsection{Offline initialization and simulator continuation}
\label{sec:method:training}

Five actors, with training seeds 6201--6205, were initialized by 50,000 IQL
updates on the immutable target replay~\citep{kostrikov2022iql}. For target
critic minimum $\bar Q_{\min}(s,a)=\min_{i\in\{1,2\}}Q_{\bar\theta_i}(s,a)$
and value function $V_{\nu}(s)$, the advantage
is $A(s,a)=\bar Q_{\min}(s,a)-V_{\nu}(s)$. Here $i$ indexes the two critic
heads, $Q_{\theta_i}$ and $Q_{\bar\theta_i}$ are the online and target
state--action value functions with parameter vectors $\theta_i$ and
$\bar\theta_i$, and $\nu$ is the value-network parameter vector.
All value functions use rewards
scaled by $\kappa_r=10^{-4}$ inside the trainer, not physical passenger-seconds.
The value uses expectile 0.7, and the
actor uses the advantage weight
\begin{equation}
w_{\mathrm{IQL}}(s,a)=
\min\left\{\exp\left(3A(s,a)\right),100\right\}.
\label{eq:iql_weight}
\end{equation}
Here $w_{\mathrm{IQL}}$ is the capped regression weight. The mixed-measure
log likelihood $\log\varpi_{\phi}(a\mid s)$ from \Cref{eq:hurdle_policy} is
used for advantage-weighted actor regression.

After a 5,000-event frozen-policy warmup, training continued in the calibrated
online simulator with 50/50 target and simulator critic minibatches and a
SAC-style actor objective~\citep{haarnoja2018soft,haarnoja2018applications}.
Define the entropy-regularized target-state bootstrap by
\begin{equation}
B_{\mathrm{SAC}}(s')=
\E_{a\sim\pi_{\phi}(\cdot\mid s')}
\left[\bar Q_{\min}(s',a)-\eta_{\mathrm{ent}}
\log\varpi_{\phi}(a\mid s')\right],
\label{eq:sac_bootstrap}
\end{equation}
where $\eta_{\mathrm{ent}}>0$ is the learned entropy temperature and the
expectation is over the hurdle policy. The discrete gate is summed exactly and
the continuous branch uses one reparameterized sample per state. The online Bellman bootstrap
mixed the fixed IQL value and entropy-regularized SAC bootstrap as
\begin{equation}
B_{\mathrm{mix}}(s')=0.1V_{\nu}(s')+0.9B_{\mathrm{SAC}}(s').
\label{eq:mixed_backup}
\end{equation}
Here $B_{\mathrm{mix}}$ is the mixed target-state value and $V_{\nu}$ is
fixed after IQL initialization. The critic target is
\begin{equation}
y_k=\kappa_r r_k+\Gamma_k B_{\mathrm{mix}}(s_{k+1}),
\label{eq:critic_target}
\end{equation}
where $y_k$ denotes a training target, distinct from the physical evaluation
endpoint $Y$. \Cref{eq:critic_target} uses no additional duration discount.
Critics minimize a Huber loss with threshold 1.0 against this target.

The actor and entropy temperature remain frozen through 10,000 simulator
events while critics adapt after warmup. On release, actor and temperature
updates use simulator states only. The actor learning-rate multiplier ramps
from approximately 0.1 to 1 over events 10,000--15,000; the temperature resumes
at its full registered rate. Each subsequent actor update is followed by
parameter interpolation toward the seed-matched 10,000-event actor, with a
fixed interpolation coefficient $5\times10^{-5}$.
The final branches inherited an exact seed-matched optimizer and replay state
at 10,000 simulator events and continued to 25,000 events. From event 15,000
through the final actor update before event 25,000, the preregistered
terminal-quarter branch linearly tapered the actor learning-rate multiplier
from 1 to 0.25 while leaving the remaining update schedule unchanged.
Each seed completed 70,000 learner updates in total. Only the 25,000-event
endpoint was exported for the final policy; checkpoints were not selected by
target-SUMO behavior.

\subsection{Dynamics-ratio diagnostic}
\label{sec:method:ratio}

The implementation retains the H2O+ domain-classifier construction for
diagnosing transition mismatch~\citep{eysenbach2021off,niu2025h2oplus}. A
joint classifier uses $(s,a,s',\Delta t/T_{\mathrm{hor}})$ and a marginal
classifier uses $(s,a)$. With balanced target and simulator classes, their
logit difference estimates the logged-target-to-simulator ratio
\begin{equation}
F_{\psi_J}(s,a,s',\Delta t/T_{\mathrm{hor}})
-G_{\psi_M}(s,a)
\approx
\log\frac{P_{\mathrm{log}}(s',\Delta t\mid s,a)}
          {P_{\mathrm{sim}}(s',\Delta t\mid s,a)}.
\label{eq:ratio_diagnostic}
\end{equation}
Here $F_{\psi_J}$ and $G_{\psi_M}$ are the joint and marginal target-class
logits, with parameter vectors $\psi_J$ and $\psi_M$, respectively; the target
class has label one and the simulator class label zero. The ratio concerns
conditional densities under a common reference measure, and its estimation
requires generalization to the states being evaluated. Because the positive
examples come from $\mathcal{M}_{\mathrm{log}}$, this estimates the logged
runtime's dynamics rather than directly estimating the formal-evaluation
kernel $P_{\mathrm{tgt}}$.
Whole trajectories, not randomly interleaved rows, define held-out validation.
The five final training runs used unit simulator weights in the critic; the
classifier output in \Cref{eq:ratio_diagnostic} remained diagnostic and never
changed a Bellman residual. This matters for interpretation: the formal effect
cannot be attributed to learned density-ratio weighting.

\subsection{Physical action calibration}
\label{sec:method:reserve}

Let $h_{\phi}(s)\in[0,60]$ be the deterministic actor proposal in seconds,
obtained from \Cref{eq:v7_action} before any operational transform, and let
$\rho=t/T_{\mathrm{hor}}$ be the horizon fraction at physical time $t$.
The first transform reduces
aggressive learned holds while retaining a transparent headway response. Let
$H_f,H_b$ be forward and backward headways, $H_f^*,H_b^*$ their targets, and
$I_f,I_b\in\{0,1\}$ indicate whether the corresponding buses are present. The
gain-one headway reserve is
\begin{equation}
h_{\mathrm{hb}}(s)=
\operatorname{clip}\left(
\tfrac{1}{2}\left[I_b(H_b-H_b^*)-I_f(H_f-H_f^*)\right],0,60
\right).
\label{eq:headway_reserve}
\end{equation}
Here $h_{\mathrm{hb}}$ is a headway-based hold in seconds, and
$\clip(\cdot,0,60)$ truncates its scalar argument to
the operational range. Headways and their targets are also measured in seconds.
The calibrated learned action combines the headway response in
\Cref{eq:headway_reserve} with the actor proposal:
\begin{equation}
h_{16}(s)=
\begin{cases}
\min\{h_{\phi}(s),\max[0.125h_{\phi}(s),h_{\mathrm{hb}}(s)]\},
& \rho<0.75\ \text{and}\ h_{\phi}(s)>0,\\
0.125h_{\phi}(s), & \text{otherwise}.
\end{cases}
\label{eq:v16_transform}
\end{equation}
Here $h_{16}$ denotes the inherited v16 calibrated holding proposal.

The completion guard is defined from physical state, not critic value. Let
$n_{\mathrm{on}}$ be onboard passengers, $c$ vehicle capacity, and
$\lambda_{\mathrm{arr}}$ the eligible passenger-arrival rate in passengers
per second. For horizon-fraction endpoints
$0\leq\rho_{\mathrm{lo}}<\rho_{\mathrm{hi}}\leq1$,
\begin{equation}
\mathcal{G}_{[\rho_{\mathrm{lo}},\rho_{\mathrm{hi}})}(s)=
\mathbb{I}\left\{
\rho_{\mathrm{lo}}\leq\rho<\rho_{\mathrm{hi}},\ I_f=1,\ I_b=0,\lambda_{\mathrm{arr}}>0,
\ c>0,\frac{n_{\mathrm{on}}}{c}\leq0.25
\right\}.
\label{eq:completion_guard}
\end{equation}
Thus the binary guard in \Cref{eq:completion_guard} identifies a low-load controlled bus with a vehicle ahead, no
vehicle behind, and continuing eligible demand.

Using the proposal in \Cref{eq:v16_transform}, the direct parent first applies an inherited 20-s tail-service floor only in
the later part of the reserve window:
\begin{equation}
h_{\mathrm{par}}(s)=
\begin{cases}
\max\{h_{16}(s),20\}, & \mathcal{G}_{[0.50,0.75)}(s)=1,\\
h_{16}(s), & \text{otherwise}.
\end{cases}
\label{eq:direct_parent}
\end{equation}
Here $h_{\mathrm{par}}$ is the direct parent's physical hold. The
reserve-augmented candidate changes only the final line of the deployed
behavior:
\begin{equation}
h_{\mathrm{safe}}(s)=
\begin{cases}
\max\{h_{\mathrm{par}}(s),60\}, &
\mathcal{G}_{[0.15,0.75)}(s)=1,\\
h_{\mathrm{par}}(s), & \text{otherwise}.
\end{cases}
\label{eq:completion_reserve}
\end{equation}
Here $h_{\mathrm{safe}}$ is the candidate's physical hold; the subscript names
the registered completion-safety variant. Because 60 s is the physical cap, an active reserve sets the hold to exactly
60 s. The candidate and direct parent share the same actor tensor, training
seed, checkpoint, state, and inherited transforms. They differ only through
\Cref{eq:completion_reserve}. The executed duration $\tau_k$ is
$h_{\mathrm{safe}}(s_k)$ for the candidate and $h_{\mathrm{par}}(s_k)$ for
its parent; the corresponding normalized action is obtained by inverting
\Cref{eq:v7_action}.

\subsection{Development and confirmation separation}
\label{sec:method:separation}

The transform was not derived from the formal outcomes. Successive
preregistered development rounds from v14 through v20 used prior development
results to diagnose actor drift, hold magnitude, headway reserve, and
completion. Several rounds reused the same nine development cells, so these
results are adaptive design evidence rather than confirmation. The 60-s
quarter-load candidate, its direct parent, all actor tensors, formal
environment blocks, primary endpoint, completion margin, and exact sign tests
were frozen before formal outcome access. The formal stage introduced no new
training run and permitted no candidate replacement.
